% Include inside a frame with % Include inside a frame with % Include inside a frame with % Include inside a frame with \input{wots_diagram}.
% H denotes one chain step; a WOTS+ step additionally uses a key and bitmask.
% \begin{tikzpicture}[
%   >=Latex,font=\small,
%   hash step/.style={draw=Teal,fill=Mist,rounded corners=1mm,
%     minimum width=7mm,minimum height=7mm,inner sep=2pt},
%   chain arrow/.style={->,Teal,line width=0.8pt}
% ]
%   \node[text=Gold] (secret) at (0,0) {$x_k$};
%   \node[hash step] (h1) at (1,0) {$H$};
%   \node (v1) at (2.4,0) {$H(x_k)$};
%   \node[hash step] (h2) at (3.9,0) {$H$};
%   \node (v2) at (5.7,0) {$H(H(x_k))$};
%   \node[hash step] (h3) at (7.5,0) {$H$};
%   \node[text=Teal,anchor=west] (public) at (8.5,0)
%     {$\mathit{pub}_k=H\bigl(H(H(x_k))\bigr)$};
%   \foreach \source/\target in {secret/h1,h1/v1,v1/h2,h2/v2,v2/h3,h3/public}
%     \draw[chain arrow] (\source) -- (\target);
%   \node[text=Muted,font=\footnotesize] at (6,-0.85)
%     {Example Winternitz chain: $w=4$, three hash steps.};
% \end{tikzpicture}
.
% H denotes one chain step; a WOTS+ step additionally uses a key and bitmask.
% \begin{tikzpicture}[
%   >=Latex,font=\small,
%   hash step/.style={draw=Teal,fill=Mist,rounded corners=1mm,
%     minimum width=7mm,minimum height=7mm,inner sep=2pt},
%   chain arrow/.style={->,Teal,line width=0.8pt}
% ]
%   \node[text=Gold] (secret) at (0,0) {$x_k$};
%   \node[hash step] (h1) at (1,0) {$H$};
%   \node (v1) at (2.4,0) {$H(x_k)$};
%   \node[hash step] (h2) at (3.9,0) {$H$};
%   \node (v2) at (5.7,0) {$H(H(x_k))$};
%   \node[hash step] (h3) at (7.5,0) {$H$};
%   \node[text=Teal,anchor=west] (public) at (8.5,0)
%     {$\mathit{pub}_k=H\bigl(H(H(x_k))\bigr)$};
%   \foreach \source/\target in {secret/h1,h1/v1,v1/h2,h2/v2,v2/h3,h3/public}
%     \draw[chain arrow] (\source) -- (\target);
%   \node[text=Muted,font=\footnotesize] at (6,-0.85)
%     {Example Winternitz chain: $w=4$, three hash steps.};
% \end{tikzpicture}
.
% H denotes one chain step; a WOTS+ step additionally uses a key and bitmask.
% \begin{tikzpicture}[
%   >=Latex,font=\small,
%   hash step/.style={draw=Teal,fill=Mist,rounded corners=1mm,
%     minimum width=7mm,minimum height=7mm,inner sep=2pt},
%   chain arrow/.style={->,Teal,line width=0.8pt}
% ]
%   \node[text=Gold] (secret) at (0,0) {$x_k$};
%   \node[hash step] (h1) at (1,0) {$H$};
%   \node (v1) at (2.4,0) {$H(x_k)$};
%   \node[hash step] (h2) at (3.9,0) {$H$};
%   \node (v2) at (5.7,0) {$H(H(x_k))$};
%   \node[hash step] (h3) at (7.5,0) {$H$};
%   \node[text=Teal,anchor=west] (public) at (8.5,0)
%     {$\mathit{pub}_k=H\bigl(H(H(x_k))\bigr)$};
%   \foreach \source/\target in {secret/h1,h1/v1,v1/h2,h2/v2,v2/h3,h3/public}
%     \draw[chain arrow] (\source) -- (\target);
%   \node[text=Muted,font=\footnotesize] at (6,-0.85)
%     {Example Winternitz chain: $w=4$, three hash steps.};
% \end{tikzpicture}
.
% H denotes one chain step; a WOTS+ step additionally uses a key and bitmask.
% \begin{tikzpicture}[
%   >=Latex,font=\small,
%   hash step/.style={draw=Teal,fill=Mist,rounded corners=1mm,
%     minimum width=7mm,minimum height=7mm,inner sep=2pt},
%   chain arrow/.style={->,Teal,line width=0.8pt}
% ]
%   \node[text=Gold] (secret) at (0,0) {$x_k$};
%   \node[hash step] (h1) at (1,0) {$H$};
%   \node (v1) at (2.4,0) {$H(x_k)$};
%   \node[hash step] (h2) at (3.9,0) {$H$};
%   \node (v2) at (5.7,0) {$H(H(x_k))$};
%   \node[hash step] (h3) at (7.5,0) {$H$};
%   \node[text=Teal,anchor=west] (public) at (8.5,0)
%     {$\mathit{pub}_k=H\bigl(H(H(x_k))\bigr)$};
%   \foreach \source/\target in {secret/h1,h1/v1,v1/h2,h2/v2,v2/h3,h3/public}
%     \draw[chain arrow] (\source) -- (\target);
%   \node[text=Muted,font=\footnotesize] at (6,-0.85)
%     {Example Winternitz chain: $w=4$, three hash steps.};
% \end{tikzpicture}
